\chapter{Résultats mesurés}
\label{ch:resultats}

Tous les chiffres de ce chapitre proviennent d'exécutions réelles sur une
installation locale (API FastAPI, frontend Vite, sources publiques interrogées en
direct). Rien n'est extrapolé : ce qui n'a pas pu être mesuré est signalé comme
tel.

\begin{encadre}[title={Environnement des captures et des mesures}]
Les captures du chapitre~\ref{ch:parcours} ont été prises dans Chrome
(1440\,\texttimes\,900) sur le frontend servi par un dev server Vite local, la
session étant ouverte avec le compte de démonstration, et les adresses de
référence diagnostiquées au préalable pour alimenter le cache local qui nourrit
le tableau de bord et le portefeuille. L'API interrogée est celle de
l'\emph{instance déployée}, pour une raison d'environnement : le poste de capture
ne joint pas \code{www.georisques.gouv.fr}. Un diagnostic y reste bloqué jusqu'au
délai de chaque appel Géorisques (mesuré : 224\,s au lieu de 10\,s) et rend les
treize aléas en \code{erreur: source Géorisques indisponible} — la capture
n'aurait montré qu'un état de panne.

Deux conséquences sont visibles sur les images et ne sont pas des défauts de
l'application : l'instance déployée est en retard sur ce dépôt, donc
\code{/api/meteo/france} y répond 404 et les cinq cartes météo de l'étape 2
affichent « Référence météo indisponible » ; et l'adresse de référence rend les
aléas selon le millésime de l'instance, non celui du poste. Ce que le contrat
exige reste vérifiable sur ces captures : une source en panne produit un
\emph{état} nommé, jamais une valeur par défaut.
\end{encadre}

\section{Le diagnostic, en conditions réelles}

Un diagnostic sur « Quai de la Rapée, 75012 Paris » a produit les faits suivants.

\begin{table}[h]
\centering
\small
\begin{tabularx}{\textwidth}{@{}p{5.2cm}X@{}}
\toprule
\textbf{Élément observé} & \textbf{Résultat} \\
\midrule
Commune retenue & Paris 12\textsuperscript{e} arrondissement, INSEE 75112 (138 024 habitants, 8 431 hab./km²) \\
Aléas du référentiel & 13 objets dans le contrat, dont 4 constatés \emph{présents} au point (ICPE, inondation, sismicité, radon) \\
Classe de crue au point & dans le périmètre TRI ; classes \emph{moyen} et \emph{faible} cartographiées au point, hauteur d'eau réglementaire disponible \\
Hydrographie & réseau reconstruit côté serveur jusqu'à la mer (Seine) \\
Météo & prévision ponctuelle de l'adresse (et non moyenne départementale) \\
Photo de terrain & Panoramax, 24 candidates, la plus proche à 18,4\,m, orientée vers le bien \\
Erreur serveur & aucune \\
\bottomrule
\end{tabularx}
\caption{Résultat d'un diagnostic réel, exécuté pendant l'écriture de ce rapport.}
\end{table}

La durée d'un diagnostic complet se situe entre 9 et 16 secondes à froid sur
l'instance déployée, la première requête payant le réveil de la fonction ; la
même adresse rejouée depuis le cache local est instantanée. Quand une source
canonique est injoignable, la durée monte au délai cumulé des appels
(224 secondes mesurées) : c'est ce cas qui a justifié la \emph{borne} côté
client, plafonnée à cinq minutes et affichée à l'écran avec le temps écoulé.
À titre de comparaison, la reconstruction manuelle des mêmes faits demande
plusieurs heures par bien.

\begin{encadre}[title={Ce que cette mesure ne dit pas}]
Elle ne dit pas que l'analyse de risque prend 15 secondes : elle dit que
l'\emph{assemblage des faits} prend 15 secondes. L'interprétation, la
tarification et la décision restent humaines — c'est précisément la répartition
que le projet revendique.
\end{encadre}

\section{Le rapport}

\begin{table}[h]
\centering
\small
\begin{tabularx}{\textwidth}{@{}p{4.6cm}X@{}}
\toprule
\textbf{Scénario} & \textbf{Mesure} \\
\midrule
Sans fournisseur de prose configuré & 200 ; \code{fallback\_used = true} ; modèle annoncé \code{template-first} ; document complet de 2 752 caractères \\
Avec prose & 200 ; \code{fallback\_used = false} ; modèle annoncé \code{gemini-2.5-flash} ; 5 narratifs de catégorie et 4 narratifs de mitigation acceptés ; 2 850 caractères \\
Rejeu de la même requête & servi par le cache (\code{cache\_hit = true}), à l'identique \\
Flux SSE & séquence complète \code{header}, \code{summary}, \code{category} \(\times\) 5, \code{mitigations}, \code{confidence}, \code{appendix}, \code{done} \\
Flux SSE avec prose & 9 événements \code{patch} de prose réellement validée \\
\bottomrule
\end{tabularx}
\caption{Le rapport, mesuré avec et sans fournisseur de prose.}
\end{table}

\section{Le fournisseur de prose}

\begin{table}[h]
\centering
\small
\begin{tabular}{@{}ll@{}}
\toprule
\textbf{Opération} & \textbf{Durée mesurée} \\
\midrule
Appel structuré (JSON) & 1,1 seconde \\
Appel en flux (SSE) & 0,7 seconde \\
\bottomrule
\end{tabular}
\caption{Latence du fournisseur de prose sur des requêtes courtes.}
\end{table}

La vérification du câblage a été faite sans clé valide, par comparaison d'erreurs :
sans en-tête d'authentification, l'API répond
« \emph{Method doesn't allow unregistered callers} » ; avec un en-tête portant une
clé invalide, elle répond \code{API\_KEY\_INVALID}. Le \emph{changement} de
réponse prouve que l'en-tête est bien celui attendu — c'est un contrôle de
câblage, pas un test de bout en bout.

\section{La suite de tests}

\begin{table}[h]
\centering
\small
\begin{tabular}{@{}lll@{}}
\toprule
\textbf{Périmètre} & \textbf{Résultat} & \textbf{Durée} \\
\midrule
Backend, suite offline & 261 tests réussis, 0 échec (2 avertissements de ressource) & 3 min 54 \\
Frontend (Vitest) & 521 tests réussis, 0 échec, 63 fichiers & 28 s \\
Lint Python & aucun avertissement & — \\
Garde-fou d'isolation du bac à sable & vert & — \\
\bottomrule
\end{tabular}
\caption{État de la suite de tests au moment de la rédaction.}
\end{table}

Un point de méthode mérite d'être souligné : la suite reste \textbf{verte alors
qu'une clé d'API réelle est configurée} dans l'environnement. C'est une propriété
importante — un test qui appellerait le modèle en direct serait lent, non
déterministe et coûteux ; ici, toutes les sources sont simulées, et la présence
d'une clé ne change rien au résultat.

\section{Les quotas, protégés côté client}

\begin{table}[h]
\centering
\small
\begin{tabularx}{\textwidth}{@{}p{3.7cm}p{3.5cm}X@{}}
\toprule
\textbf{Source} & \textbf{Plafond appliqué} & \textbf{Comportement en cas de refus} \\
\midrule
BDNB & 120 requêtes / minute & budget mensuel de 10 000 appels ; retour partiel déclaré \\
Géocodage (Géoplateforme) & 50 requêtes / seconde & une seule réattente sur un 429 explicite, puis échec franc \\
Routes de collecte & 30 requêtes / 60 s & réponse 429 côté API, jamais un ralentissement silencieux \\
\bottomrule
\end{tabularx}
\caption{Les plafonds vérifiés chez le fournisseur et leur application.}
\end{table}
